\origpage{54}{N. Robertson, P.D. Seymour / Journal of Combinatorial Theory, Series B 89 (2003) 43--76}
Now $N_1$ has maximum degree $\le 3$, and there is a subgraph $N_4$ of $G$ isomorphic to a subdivision of $N_1$ (and hence $N_1$ is a minor of $N_4$). We may assume that $N_4$ is a rigid drawing in a sphere $\Sigma$. Let $\mathscr{T}_4$ be the tangle in $N_4$ induced by $\mathscr{T}_1$. Then $\mathscr{T}_3$ is the tangle in $G$ induced by $\mathscr{T}_4$, and since $\mathscr{T}_3 \subseteq \mathscr{T}^*$ it follows that $N_4, \mathscr{T}_4$ is a $\Sigma$-span of order $\theta$. This proves 4.2. $\square$

Our second lemma is the following, theorem 4.3 of [7].

\textbf{4.3.} \textit{If $L$ is a graph and $\Sigma$ is a surface such that $L$ can be drawn in $\Sigma$, there is an integer $\theta \ge 1$ with the following property. If $\mathscr{T}^*$ is a tangle in a graph $G$ controlling no $L$ minor, then there is no $\Sigma$-span of order $\ge \theta$.}

\textbf{Proof of 3.1} (assuming 4.1). Up to homeomorphism, let $\Sigma_1, \dots, \Sigma_n$ be all the surfaces on which $L$ cannot be drawn together with the sphere and all the surfaces which can be obtained from a surface on which $L$ cannot be drawn by adding a handle or crosscap. (There are only finitely many up to homeomorphism.) Let us number $\Sigma_1, \dots, \Sigma_n$ so that for $1 \le i, j \le n$, if $\Sigma_j$ can be obtained from $\Sigma_i$ by adding a handle or crosscap then $j \ge i$.

Now we define $\kappa_i, \rho_i, \zeta_i$ for $1 \le i \le n$ as follows. Choose $p \ge 0$ so that $K_p$ has a minor isomorphic to $L$. Inductively, suppose that $1 \le i \le n$, and $\kappa_j, \rho_j, \zeta_j, \theta_j$ are defined for all $j$ with $i \delm{\le}\fix{<} j \le n$.\fixnote{fix:3.1-proof}{Here $\theta_i$ is the one being chosen, so only $j > i$ is assumed.} If $L$ can be drawn in $\Sigma_i$, choose $\theta_i$ so that 4.3 is satisfied (with $\Sigma, \theta$ replaced by $\Sigma_i, \theta_i$) and let $\kappa_i, \rho_i, \zeta_i = 0$. If $L$ cannot be drawn in $\Sigma_i$, then $i < n$; choose $\theta_i, \kappa_i, \rho_i, \zeta_i$ so that 4.1 is satisfied (with $\phi, \kappa, \rho, \zeta, \theta$ replaced by $\max(\theta_{i+1}, \dots, \theta_n), \kappa_i, \rho_i, \zeta_i, \theta_i$). This completes the inductive definition.

Choose $\theta^* \ge 1$ so that 4.2 is satisfied (with $\theta$ replaced by $\theta_1$). Let $\kappa = \max(\kappa_1, \dots, \kappa_n)$, $\rho = \max(\rho_1, \dots, \rho_n)$, $\zeta = \max(\zeta_1, \dots, \zeta_n)$ and $\theta = \theta^*$. We claim that $\kappa, \rho, \zeta, \theta$ satisfy 3.1.

For let $\mathscr{T}^*$ be a tangle of order $\ge \theta = \theta^*$ in a graph $G$, controlling no $L$-minor. By 4.2, there is a $\Sigma_1$-span of order $\ge \theta_1$ (because $\Sigma_1$ is the sphere). Choose $i$ with $1 \le i \le n$ maximum so that there is a $\Sigma_i$-span of order $\ge \theta_i$. If $L$ can be drawn in $\Sigma_i$ then $\theta_i$ satisfies 4.3 and yet $\mathscr{T}^*$ controls no $L$-minor, a contradiction. Thus, $L$ cannot be drawn in $\Sigma_i$, and so one of the outcomes of 4.1 holds (with $\phi, \kappa, \rho, \zeta, \theta$ replaced by $\max(\theta_{i+1}, \dots, \theta_n), \kappa_i, \rho_i, \zeta_i, \theta_i$). Since $L$ cannot be drawn in $\Sigma_i$ it follows that every surface which can be obtained from $\Sigma_i$ by adding a handle or crosscap occurs in the list $\Sigma_{i+1}, \dots, \Sigma_n$ (up to homeomorphism), and so the first outcome of 4.1 does not hold from the maximality of $i$. Moreover the third outcome of 4.1 does not hold, since $\mathscr{T}^*$ controls no $L$-minor and hence controls no $K_p$-minor, by theorem 4.2 of [7]. Thus the second outcome of 4.1 holds, and so 3.1 holds as required. $\square$

\sectionhead{5. Induction on horns}

Thus, it remains to prove 4.1. Let $\Sigma$ be a surface, let $\mathscr{T}^*$ be a tangle in a graph $G$, and let $H, \eta, \mathscr{T}$ be a $\Sigma$-span in $G$ with respect to $\mathscr{T}^*$. Let $v \in V(G) \setminus V(\eta(H))$.
